Examples drawn from the MPI Standard document

A simple `hello world' example usage of point-to-point communication.

%% FIXME: Use MPI_STATUS_IGNORE instead of declare, ignore status?
%%HEADER
%%LANG: C
%%DECL:#include <string.h>
%%ENDHEADER
\begin{lstlisting}[language={[MPI]C}]
#include "mpi.h"
int main(int argc, char *argv[])
{
  char message[20];
  int myrank;
  MPI_Status status;
  MPI_Init(&argc, &argv);
  MPI_Comm_rank(MPI_COMM_WORLD, &myrank);
  if (myrank == 0)    /* code for process zero */
  {
      strcpy(message,"Hello, there");
      MPI_Send(message, strlen(message)+1, MPI_CHAR, 1, 99,
               MPI_COMM_WORLD);
  }
  else if (myrank == 1)  /* code for process one */
  {
      MPI_Recv(message, 20, MPI_CHAR, 0, 99, MPI_COMM_WORLD, &status);
      printf("received :%s:\n", message);
  }
  MPI_Finalize();
  return 0;
}
\end{lstlisting}
\end{example}

In Example~\ref{ex:pt2pt-helloworld}, process zero (\code{myrank = 0}) sends a \mpiterm{message} to process one
using the
\mpiterm{send} operation \mpifunc{MPI\_SEND}.

The following examples illustrate the first two cases.
\label{page:pt2pt-3examples}

\begin{example}
\label{pt2pt-exA}
\exindex{MPI\_SEND}%
\exindex{MPI\_RECV}%
\exindex{Datatypes!matching}%
Sender and receiver specify matching types.
%%HEADER
%%LANG: FORTRAN
%%FRAGMENT
%%DECL: integer comm, rank, ierr, tag, a(2), b(2)
%%DECL: integer status(MPI_STATUS_SIZE)
%%ENDHEADER
\begin{lstlisting}[language={[MPI]Fortran}]
CALL MPI_COMM_RANK(comm, rank, ierr)
IF (rank .EQ. 0) THEN
   CALL MPI_SEND(a(1), 10, MPI_REAL, 1, tag, comm, ierr)
ELSE IF (rank .EQ. 1) THEN
   CALL MPI_RECV(b(1), 15, MPI_REAL, 0, tag, comm, status, ierr)
END IF
\end{lstlisting}

\begin{example}
\label{pt2pt-exB}
\exindex{MPI\_SEND}%
\exindex{MPI\_RECV}%
\exindex{Datatypes!not matching}%
Sender and receiver do not specify matching types.

   CALL MPI_RECV(recvbuf, count, MPI_REAL, 1, tag, comm, status, ierr)
ELSE IF (rank .EQ. 1) THEN
   CALL MPI_SEND(sendbuf, count, MPI_REAL, 0, tag, comm, ierr)
   CALL MPI_RECV(recvbuf, count, MPI_REAL, 0, tag, comm, status, ierr)
END IF
\end{lstlisting}
The message sent by each process has to be copied out before the send operation
returns and the receive operation starts.  For the program to complete, it is
necessary that at least one of the two messages sent be buffered.
Thus, this program can
succeed only if
the communication system can buffer at least \mpiarg{count} words of data.
\end{example}


\begin{example}
\label{pt2pt-exZ}
\exindex{MPI\_Buffer\_attach}%
\exindex{MPI\_Buffer\_detach}%
Calls to attach and detach buffers.
%%HEADER
%%LANG: C
%%FRAGMENT
%%ENDHEADER
\begin{lstlisting}[language={[MPI]C}]
#define BUFFSIZE 10000
int size;
char *buff;
MPI_Buffer_attach(malloc(BUFFSIZE), BUFFSIZE);
/* a buffer of 10000 bytes can now be used by MPI_Bsend */
MPI_Buffer_detach(&buff, &size);
/* Buffer size reduced to zero */
MPI_Buffer_attach(buff, size);
/* Buffer of 10000 bytes available again */
\end{lstlisting}
\end{example}



\begin{example}
\label{pt2pt-exJ}
\exindex{MPI\_ISEND}%
\exindex{MPI\_IRECV}%
\exindex{MPI\_WAIT}%
\exindex{Nonblocking operations}%
Simple usage of nonblocking operations and \mpifunc{MPI\_WAIT}.
%%HEADER
%%LANG: FORTRAN
%%FRAGMENT
%%DECL: integer comm, rank, ierr, tag, request
%%DECL: real a(10)
%%DECL: integer status(MPI_STATUS_SIZE)
%%SUBST:\*\*.*:tag=1
%%ENDHEADER
\begin{lstlisting}[language={[MPI]Fortran}]
CALL MPI_COMM_RANK(comm, rank, ierr)
IF (rank .EQ. 0) THEN
   CALL MPI_ISEND(a(1), 10, MPI_REAL, 1, tag, comm, request, ierr)
   ! **** do some computation to mask latency ****
   CALL MPI_WAIT(request, status, ierr)
ELSE IF (rank .EQ. 1) THEN
   CALL MPI_IRECV(a(1), 15, MPI_REAL, 0, tag, comm, request, ierr)
   ! **** do some computation to mask latency ****
   CALL MPI_WAIT(request, status, ierr)
END IF
\end{lstlisting}
\end{example}

A request object can be \mpitermni{freed}\mpitermindex{request objects!freeing} using the following \MPI/ procedure.

\begin{example}
\label{pt2pt-exK}
\exindex{MPI\_REQUEST\_FREE}%
\exindex{MPI\_ISEND}%
\exindex{MPI\_IRECV}%
\exindex{MPI\_WAIT}%
\exindex{Nonblocking operations}%
An example using \mpifunc{MPI\_REQUEST\_FREE}.

%%HEADER
%%LANG: FORTRAN
%%FRAGMENT
%%DECL: integer rank, req, ierr, outval, inval, i, n
%%DECL: integer status(MPI_STATUS_SIZE)
%%ENDHEADER
\begin{lstlisting}[language={[MPI]Fortran},basicstyle=\lstsmall]
CALL MPI_COMM_RANK(MPI_COMM_WORLD, rank, ierr)
IF (rank .EQ. 0) THEN
   DO i=1,n
      CALL MPI_ISEND(outval, 1, MPI_REAL, 1, 0, MPI_COMM_WORLD, req, ierr)
      CALL MPI_REQUEST_FREE(req, ierr)
      CALL MPI_IRECV(inval, 1, MPI_REAL, 1, 0, MPI_COMM_WORLD, req, ierr)
      CALL MPI_WAIT(req, status, ierr)
   END DO
ELSE IF (rank .EQ. 1) THEN
   CALL MPI_IRECV(inval, 1, MPI_REAL, 0, 0, MPI_COMM_WORLD, req, ierr)
   CALL MPI_WAIT(req, status, ierr)
   DO I=1,n-1
      CALL MPI_ISEND(outval, 1, MPI_REAL, 0, 0, MPI_COMM_WORLD, req, ierr)
      CALL MPI_REQUEST_FREE(req, ierr)
      CALL MPI_IRECV(inval, 1, MPI_REAL, 0, 0, MPI_COMM_WORLD, req, ierr)
      CALL MPI_WAIT(req, status, ierr)
   END DO
   CALL MPI_ISEND(outval, 1, MPI_REAL, 0, 0, MPI_COMM_WORLD, req, ierr)
   CALL MPI_WAIT(req, status, ierr)
END IF
\end{lstlisting}
\end{example}

\subsection{Semantics of Nonblocking Communications}
\mpitermtitleindex{semantics!nonblocking communications}
\label{subsec:pt2pt-semantics}

\begin{example}
\label{coll-exX-fortran}
\exindex{MPI\_Reduce}%
\exindex{MPI\_Op\_create}%
\exindex{MPI\_User\_function}%
\mpicallbackindex{MPI\_User\_function}%
How to use the \ftype{mpi\_f08} interface of the Fortran \ftype{MPI\_User\_function}.
%%HEADER
%%LANG: FORTRAN08
%%ENDHEADER
\begin{lstlisting}[language={[MPI08]Fortran}]
subroutine my_user_function(invec, inoutvec, len, dtype)   bind(c)
   use, intrinsic :: iso_c_binding, only : c_ptr, c_f_pointer
   use mpi_f08
   type(c_ptr), value :: invec, inoutvec
   integer :: len
   type(MPI_Datatype) :: dtype
   real, pointer :: invec_r(:), inoutvec_r(:)
   if (dtype%MPI_VAL == MPI_REAL%MPI_VAL) then
      call c_f_pointer(invec, invec_r, (/ len /))
      call c_f_pointer(inoutvec, inoutvec_r, (/ len /))
      inoutvec_r = invec_r + inoutvec_r
   end if
end subroutine
\end{lstlisting}
\end{example}

%%HEADER
%%LANG: FORTRAN90
%%ENDHEADER
\begin{lstlisting}[language={[MPI]Fortran}]
SUBROUTINE SUM(A, B, map, m, comm, p)
USE MPI
INTEGER m, map(m), comm, p, win, ierr, disp_int, i, j
REAL A(m), B(m)
INTEGER(KIND=MPI_ADDRESS_KIND) lowerbound, size, realextent, disp_aint

CALL MPI_TYPE_GET_EXTENT(MPI_REAL, lowerbound, realextent, ierr)
size = m * realextent
disp_int = realextent
CALL MPI_WIN_CREATE(B, size, disp_int, MPI_INFO_NULL,  &
                    comm, win, ierr)

CALL MPI_WIN_FENCE(0, win, ierr)
DO i=1,m
   j = map(i)/m
   disp_aint = MOD(map(i),m)
   CALL MPI_ACCUMULATE(A(i), 1, MPI_REAL, j, disp_aint, 1, MPI_REAL,   &
                       MPI_SUM, win, ierr)
END DO
CALL MPI_WIN_FENCE(0, win, ierr)

CALL MPI_WIN_FREE(win, ierr)
RETURN
END
\end{lstlisting}

%%HEADER
%%LANG: C
%%FRAGMENT
%%SKIPELIPSIS
%%DECL: int converged(int), i, A; MPI_Win win;
%%DECL: int toneighbors, toneighbor[10], frombuf[10]; int update(int);
%%DECL: MPI_Datatype totype[10], fromtype[10];
%%DECL: MPI_Aint todisp[10];
%%ENDHEADER
\begin{lstlisting}[language={[MPI]C}]
...
while (!converged(A)) {
  update(A);
  MPI_Win_fence(MPI_MODE_NOPRECEDE, win);
  for(i=0; i < toneighbors; i++)
    MPI_Put(&frombuf[i], 1, fromtype[i], toneighbor[i],
                         todisp[i], 1, totype[i], win);
  MPI_Win_fence((MPI_MODE_NOSTORE | MPI_MODE_NOSUCCEED), win);
}
\end{lstlisting}

\exindex{MPI\_TYPE\_CREATE\_SUBARRAY}%
%%HEADER
%%LANG: C
%%FRAGMENT
%%DECL:void foobar(double [100][25]);
%%TAIL:(void)foobar(subarray);
%%ENDHEADER
\begin{lstlisting}[language={[MPI]C}]
   double subarray[100][25];
   MPI_Datatype filetype;
   int sizes[2], subsizes[2], starts[2];
   int rank;

   MPI_Comm_rank(MPI_COMM_WORLD, &rank);
   sizes[0]=100; sizes[1]=100;
   subsizes[0]=100; subsizes[1]=25;
   starts[0]=0; starts[1]=rank*subsizes[1];

   MPI_Type_create_subarray(2, sizes, subsizes, starts, MPI_ORDER_C,
                            MPI_DOUBLE, &filetype);
\end{lstlisting}

Or, equivalently in Fortran:

%%HEADER
%%LANG: FORTRAN
%%FRAGMENT
%%TAIL: call foobar(subarray)
%%ENDHEADER
\begin{lstlisting}[language={[MPI]Fortran}]
double precision subarray(100,25)
integer filetype, rank, ierror
integer sizes(2), subsizes(2), starts(2)

call MPI_COMM_RANK(MPI_COMM_WORLD, rank, ierror)
sizes(1)    = 100
sizes(2)    = 100
subsizes(1) = 100
subsizes(2) = 25
starts(1)   = 0
starts(2)   = rank*subsizes(2)

call MPI_TYPE_CREATE_SUBARRAY(2, sizes, subsizes, starts, &
           MPI_ORDER_FORTRAN, MPI_DOUBLE_PRECISION,       &
           filetype, ierror)
\end{lstlisting}

%%HEADER
%%LANG: C
%%SKIPELIPSIS
%%TOP:#define MAX_DATA 10
%%TOP:#define FATAL 12
%%TOP:void error(int,char *);
%%ENDHEADER
\begin{lstlisting}[language={[MPI]C}]
#include "mpi.h"
int main(int argc, char *argv[])
{
    MPI_Comm client;
    MPI_Status status;
    char port_name[MPI_MAX_PORT_NAME];
    double buf[MAX_DATA];
    int    size, again;

    MPI_Init(&argc, &argv);
    MPI_Comm_size(MPI_COMM_WORLD, &size);
    if (size != 1) error(FATAL, "Server too big");
    MPI_Open_port(MPI_INFO_NULL, port_name);
    printf("server available at %s\n", port_name);
    while (1) {
        MPI_Comm_accept(port_name, MPI_INFO_NULL, 0, MPI_COMM_WORLD, 
                        &client);
        again = 1;
        while (again) {
            MPI_Recv(buf, MAX_DATA, MPI_DOUBLE, 
                     MPI_ANY_SOURCE, MPI_ANY_TAG, client, &status);
            switch (status.MPI_TAG) {
                case 0: MPI_Comm_free(&client);
                        MPI_Close_port(port_name);
                        MPI_Finalize();
                        return 0;
                case 1: MPI_Comm_disconnect(&client);
                        again = 0;
                        break;
                case 2: /* do something */
                ...
                default:
                        /* Unexpected message type */
                        MPI_Abort(MPI_COMM_WORLD, 1);
                }
            }
        }
}
\end{lstlisting}

%%HEADER
%%SKIP
%%ENDHEADER
\begin{verbatim}
    mpiexec spec0 [: spec1 : spec2 : ...]
\end{verbatim}
\const{MPI\_APPNUM} should be set to the number of the corresponding
specification.


The following code fragment shows
some possible ways to send scalars or arrays of interoperable
derived types in Fortran. The example assumes that all data is passed by address.
\exindex{Fortran 90!derived types}%
\exindex{MPI\_TYPE\_CREATE\_RESIZED}%
\exindex{MPI\_TYPE\_CREATE\_STRUCT}%
\exindex{MPI\_GET\_ADDRESS}%
\exindex{MPI\_TYPE\_COMMIT}%
%%HEADER
%%LANG: FORTRAN
%%EXTRAFLAGS:mismatch
%%FRAGMENT
%%DECL: integer ierr, dest, tag, comm
%%DECL: integer newtype, newarrtype
%%ENDHEADER
\begin{lstlisting}[language={[MPI]Fortran}]
type, BIND(C) :: mytype
   integer :: i
   real :: x
   double precision :: d
   logical :: l
end type mytype

type(mytype) :: foo, fooarr(5)
integer :: blocklen(4), dtype(4)
integer(KIND=MPI_ADDRESS_KIND) :: disp(4), base, lb, extent

call MPI_GET_ADDRESS(foo%i, disp(1), ierr)
call MPI_GET_ADDRESS(foo%x, disp(2), ierr)
call MPI_GET_ADDRESS(foo%d, disp(3), ierr)
call MPI_GET_ADDRESS(foo%l, disp(4), ierr)

base = disp(1)
disp(1) = disp(1) - base
disp(2) = disp(2) - base
disp(3) = disp(3) - base
disp(4) = disp(4) - base

blocklen(1) = 1
blocklen(2) = 1
blocklen(3) = 1
blocklen(4) = 1

dtype(1) = MPI_INTEGER
dtype(2) = MPI_REAL
dtype(3) = MPI_DOUBLE_PRECISION
dtype(4) = MPI_LOGICAL

call MPI_TYPE_CREATE_STRUCT(4, blocklen, disp, dtype, newtype, ierr)
call MPI_TYPE_COMMIT(newtype, ierr)

call MPI_SEND(foo%i, 1, newtype, dest, tag, comm, ierr)
! or
call MPI_SEND(foo, 1, newtype, dest, tag, comm, ierr)
! expects that base == address(foo%i) == address(foo)

call MPI_GET_ADDRESS(fooarr(1), disp(1), ierr)
call MPI_GET_ADDRESS(fooarr(2), disp(2), ierr)
extent = disp(2) - disp(1)
lb = 0
call MPI_TYPE_CREATE_RESIZED(newtype, lb, extent, newarrtype, ierr)
call MPI_TYPE_COMMIT(newarrtype, ierr)

call MPI_SEND(fooarr, 5, newarrtype, dest, tag, comm, ierr)
\end{lstlisting}
\exindex{MPI\_TYPE\_CREATE\_RESIZED}%
\exindex{MPI\_TYPE\_CREATE\_STRUCT}%
\exindex{MPI\_GET\_ADDRESS}%
\exindex{MPI\_TYPE\_COMMIT}%

%%HEADER
%%LANG: C
%%FRAGMENT
%%DECL: MPI_Comm comm; int ndims, d, rank_source, rank_dest;
%%DECL: int sendbuf[1], sendcount, recvbuf[1], recvcount;
%%DECL: MPI_Datatype recvtype, sendtype;
%%DECL: MPI_Status status;
%%ENDHEADER
\begin{lstlisting}[language={[MPI]C},basicstyle=\lstsmall,escapeinside=`']
MPI_Cartdim_get(comm, &ndims);
for( /*direction*/ d=0; d < ndims; d++) {
    MPI_Cart_shift(comm, /*direction*/ d, /*disp*/ 1, &rank_source, &rank_dest);
    MPI_Sendrecv(&sendbuf[d*2`\underline{+0}'],sendcount,sendtype,`\underline{rank\_source}',/*sendtag*/d*2,
                 &recvbuf[d*2`\underline{+1}'],recvcount,recvtype,`\underline{rank\_dest}', /*recvtag*/ d*2,
                 comm,&status);/*communication in direction of displacment -1*/
    MPI_Sendrecv(&sendbuf[d*2`\underline{+1}'],sendcount,sendtype,`\underline{rank\_dest}', /*sendtag*/ d*2+1,
                 &recvbuf[d*2`\underline{+0}'],recvcount,recvtype,`\underline{rank\_source}',/*recvtag*/d*2+1,
                 comm,&status);/*communication in direction of displacment +1*/
}
\end{lstlisting}

